\subsection{The dyadic Prym map and its exact torsion fibres}

\begin{statement}{A.152}{Proved}
{An injective dyadic Kummer map and the boundary torsion test}
Retain \(K=\Q(\sqrt2)\), \(\epsilon=1+\sqrt2\),
\(M=325596479663\), and \(E_b/K\) from A.150.
Write \(K_2=\Q_2(\sqrt2)\), \(\pi=\sqrt2\),
and normalize \(v_\pi(\pi)=1\).
For every integer \(n\geq1\), the actual multiplication map
on this elliptic curve satisfies
\[
 [2^n]E_b(K_2)\ \subseteq\
 \{O\}\ \cup\
 \{(x,y)\in E_b(K_2):v_\pi(x)\leq2-4n\}.
\tag{A.152.1}
\]
Consequently the continuous Kummer homomorphism
\[
 \kappa_2:E_b(K_2)\longrightarrow
                   H^1_{\rm cts}(G_{K_2},T_2E_b)
\tag{A.152.2}
\]
is injective, and the exact local and global torsion subgroups are
\[
             E_b(K_2)_{\rm tors}=E_b(K)_{\rm tors}=\{O,(0,0)\}.
\tag{A.152.3}
\]
Let \(q_+:\overline D_K\to E_b\) be the positive-sign
degree-two morphism of A.150, including its displayed integral
scaling.  Denote the four boundary points of \(\overline D\)
by \(P_0,P_1,P_{1/2},P_\infty\), according to their
\(T\)-coordinate.  There is an exact equality of subsets
of \(\overline D(K_2)\):
\[
 (\kappa_2\circ q_+)^{-1}
       \bigl(\kappa_2(E_b(K)_{\rm tors})\bigr)
                =\{P_0,P_1,P_{1/2},P_\infty\}.
\tag{A.152.4}
\]
On the left, the global points are first included into
\(E_b(K_2)\).

For the section comparison, let \(\mathscr S_2\) be the set
of continuous section classes of \(\pi_1(D^\circ/\Q)\to G_\Q\)
whose localization at two is an ordinary point section of
\(D^\circ(\Q_2)\).  Restriction to \(G_K\), compactification,
and \(q_+\), followed by the geometric two-primary quotient
of the elliptic fundamental group, define the actual map
\[
 \eta_2:\mathscr S_2\longrightarrow H^1_{\rm cts}(G_K,T_2E_b).
\tag{A.152.5}
\]
Its image satisfies
\[
          \eta_2(\mathscr S_2)\ \cap\
                \kappa_2(E_b(K)_{\rm tors})=\varnothing .
\tag{A.152.6}
\]
This includes full Selmer sections ordinary at two.
The theorem makes no conclusion about a section from a
single cuspidal localization, and uses no local section theorem.
\end{statement}
\begin{proof}
\emph{The literal doubling calculation.}
The coefficient has the exact form
\[
             b=\pi^2\,573(\epsilon M)^2,
             \qquad v_\pi(b)=2,\qquad 573\equiv5\pmod8.
\tag{A.152.7}
\]
The tangent formula on \(y^2=x^3+bx\) gives, for a point
outside the kernel of doubling,
\[
                   x([2]P)=\frac{(x^2-b)^2}{4x(x^2+b)}.
\tag{A.152.8}
\]
We determine its valuation for every possible value
\(r=v_\pi(x)\).
For \(r<0\), the smaller valuation in \(x^2\pm b\)
is \(2r\), so (A.152.8) has valuation \(r-4\).
For \(r=0\) its valuation is \(-4\).
For \(r>1\), the smaller valuation in both terms is two,
and the resulting valuation is \(-r-2\).

It remains to handle \(r=1\), where those two valuations
coincide.  Write \(x=\pi\epsilon Mz\) with \(z\) a unit.
The unit squares in
\(\mathcal O_{K_2}/\pi^5\), written as
\(a+c\pi\) with \(a\bmod8\) and \(c\bmod4\), are
exactly
\[
                                1,\qquad3+2\pi.
\tag{A.152.9}
\]
For completeness, squaring the sixteen units
\(a+c\pi\), with \(a\) odd, gives
\((a^2+2c^2)+2ac\pi\); its reduction is one of the two
displayed values, and both occur.
The first value gives
\(v_\pi(z^2+573)=2\), so the right side of the curve
equation has valuation
\[
              v_\pi\bigl(x(x^2+b)\bigr)=1+2+2=5.
\]
It therefore cannot equal \(y^2\).
For the second value, the exact residues give
\[
             v_\pi(z^2+573)=3,\qquad
             v_\pi(z^2-573)=2.
\tag{A.152.10}
\]
Thus \(v_\pi(x^2+b)=5\), \(v_\pi(x^2-b)=4\),
and (A.152.8) has valuation \(8-(4+1+5)=-2\).
These same residues show that \(x^2=-b\) has no solution
in \(K_2\): such a solution would have \(r=1\), whereas
both possible residues of \(z^2+573\) are nonzero.
The kernel of doubling is consequently \(\{O,(0,0)\}\).

We have proved the exhaustive table
\[
 \begin{array}{c|c}
 v_\pi(x(P))&v_\pi(x([2]P))\\\hline
 r<0&r-4\\
 0&-4\\
 1&-2\\
 r>1&-r-2.
 \end{array}
\tag{A.152.11}
\]
In particular every nonzero double has valuation at most
minus two.  Every further nonzero double reduces that
valuation by exactly four.  Induction proves (A.152.1),
including every point whose intermediate double is the origin.
It follows directly that
\[
                        \bigcap_{n\geq1}[2^n]E_b(K_2)=\{O\}:
\tag{A.152.12}
\]
the \(n=1\) inclusion gives \(x\ne0\) for every nonzero
member of this intersection.  Its finite \(x\)-valuation
cannot satisfy the bounds in (A.152.1) for all \(n\).

An odd-order torsion point belongs to every image
\([2^n]E_b(K_2)\), because multiplication by \(2^n\)
is invertible on its finite cyclic subgroup.  Equation
(A.152.12) excludes every such nonzero point.
A point of two-power order at least four has a multiple
whose double is the unique nonzero two-torsion point
\((0,0)\).  The \(n=1\) inclusion excludes that double.
Decomposing a finite-order cyclic group into its odd and
two-primary parts proves \(E_b(K_2)_{\rm tors}=\{O,(0,0)\}\).
The field inclusion \(K\hookrightarrow K_2\) gives an
injective map on curve points and preserves the two
displayed torsion points, proving the global equality
in (A.152.3) as well.

\emph{The continuous Kummer map.}
For \(P\in E_b(K_2)\), choose compatible division points
\(P_n\in E_b(\overline K_2)\) with \(P_0=P\) and
\([2]P_{n+1}=P_n\).
The cocycle
\[
            g\longmapsto(gP_n-P_n)_{n\geq1}\in T_2E_b
\tag{A.152.13}
\]
is continuous because each finite-coordinate cocycle is
continuous.  A different compatible choice changes it by
a coboundary in \(T_2E_b\), so it defines (A.152.2).
Its projection to \(E_b[2^n]\) is the connecting class of
the exact multiplication sequence.
Vanishing of (A.152.13) therefore puts \(P\) in
\([2^n]E_b(K_2)\) for every \(n\).
Equation (A.152.12) proves injectivity, without replacing
continuous cohomology by an unjustified inverse limit.

\emph{An independent finite-reduction torsion bound.}
The two primes seventeen and forty-one split in \(K\).
For each chosen root \(t_0\) of \(t^2=2\) in the
indicated residue field, the residue homomorphism is
\(\sqrt2\mapsto t_0\).
The complete point counts on the displayed smooth cubic
models are
\[
 \begin{array}{c|c|c|c}
 p&t_0&b\bmod p&\#E_b(\F_p)\\\hline
 17&6&10&10\\
 17&11&11&10\\
 41&17&21&52\\
 41&24&9&52.
 \end{array}
\tag{A.152.14}
\]
Every entry is checked by enumerating all pairs in
\(\F_p^2\) and then adding the origin.  The discriminant
\(-64b^3\) is nonzero at all four reductions.
The prime-to-\(p\) torsion injects under the actual
reduction homomorphism
\cite[II, Corollary 4.2(b) and Remark 4.3]{AgentMilneEC2021}.
For a coefficient prime different from seventeen and
forty-one, its torsion order therefore divides both ten
and fifty-two.
The seventeen-primary torsion is zero by reduction at
forty-one, and the forty-one-primary torsion is zero by
reduction at seventeen.
Thus the total torsion has order dividing two.
The point \((0,0)\) has order two, proving (A.152.3).

\emph{The proper fibres and the arithmetic section map.}
Before the integral scaling of A.150 the quotient functions
are
\[
              x_+=\frac{aW^2}{T},\qquad
              y_+=\frac{aW(T+1/\sqrt2)}{T}.
\tag{A.152.15}
\]
At \(P_0\) and \(P_\infty\) their pole orders are
two and three, and at \(P_1,P_{1/2}\) both functions
vanish.  Each of the two pairs is exchanged by the
involution of A.150.  The morphism has degree two and
is unramified at these four points.  Therefore its
scheme-theoretic fibres are exactly
\[
                q_+^{-1}(O)=P_0+P_\infty,
              \qquad q_+^{-1}((0,0))=P_1+P_{1/2}.
\tag{A.152.16}
\]
Scaling preserves both target points.
Combining these exact fibres with the injectivity in
(A.152.2) proves (A.152.4).

Here are all comparison maps used for (A.152.5).
Base change gives
\(\pi_1(D^\circ_K)=\pi_1(D^\circ)\times_{G_\Q}G_K\),
so restriction of a section gives a section of this
fibre product.
The open immersion into \(\overline D_K\) induces
the arithmetic fundamental-group map killing boundary
inertia.  The proper morphism \(q_+\) induces the next
arithmetic map to \(\pi_1(E_b)\).
Finally the multiplication covers \([2^n]:E_b\to E_b\)
give its geometric two-primary quotient
\[
                  T_2E_b\rtimes G_K\longrightarrow G_K,
\tag{A.152.17}
\]
with splitting supplied by the origin.
These are the actual covers: over the algebraic closure,
a connected finite etale cover of an elliptic curve has
genus one by Riemann--Hurwitz.  Choosing an origin over
the origin makes it an isogeny, and the dual isogeny
shows that it is dominated by a multiplication cover.
Thus the stated quotient is indeed the geometric
two-primary quotient, with its standard Tate-module action.
Conjugation by a geometric element changes the associated
continuous cocycle by a coboundary, so (A.152.5) is defined
on the stated section classes.

Every map just described commutes with restriction to
the decomposition group of the unique place of \(K\)
over two.  An ordinary point \(P\in D^\circ(\Q_2)\)
base changes to \(D^\circ(K_2)\); its image in
(A.152.17) is precisely the cocycle (A.152.13) for
\(q_+(P)\).
Equation (A.152.4) excludes either global torsion
Kummer class at this localization.  A global equality
would restrict to that excluded local equality.
This proves (A.152.6).
\end{proof}

The executed certificate
\path{code/harmonic_prym_dyadic_injection_certificate.py}
retains the complete unit-square enumeration, the exact
duplication identity, the valuation-case identities, and
all four finite-field point lists.
The separate retained Magma local calculation reports a
minimal model with Kodaira type \(I_4^*\), minimal
discriminant valuation eighteen, and component order four.
Those software outputs are a comparison check; the proof
of (A.152.1)--(A.152.12) is the explicit valuation argument.
